\honest{Purchase Fuzzies and Utilons Separately}{Eliezer Yudkowsky}{2009}

I begin by quoting yesterday's post: the lawyer at the soup kitchen, who may volunteer to stay motivated if he also gives money.

I hold open doors for little old ladies. Recently, on a walk, I saw a station wagon with its trunk wide open, checked that no one was unloading it, and knocked to tell the owners. People do such things from concern, fear of guilt, a wish to signal trustworthiness, or because helping is pleasant. All of these motives are fine by me; I might give bonus points for the first, but deduct none for the others, ``Just so long as people get helped.'' \nb{Economists call the last motive the ``warm glow'' of giving, after James Andreoni's models of 1989 and 1990.}

But I work in the nonprofit sector. Was that minute the best use of my time? The obvious defense, ``or perhaps, obvious rationalization,'' is that the deed restores my willpower better than listening to music. I also suspect that if I walk past problems, my altruism will fade. I have not tested this. \nb{The willpower link is to ego depletion. In 2016 a replication in 23 laboratories found that effect close to zero. A 2010 study found that people held a weight longer after donating to charity; studies of ``moral licensing'' suggest that a good deed can make the next one less likely in the short term. None of this settles the author's case.}

But then the benefits I list are to myself. Who said I was defending it as a selfless deed? It is a ``selfish good deed,'' with indirect benefits to others if it keeps me altruistic. If you distrust it as an ulterior motive, you could as well look straight at the good deed behind it. \nb{Yudkowsky had written in ``Circular Altruism'' (2008) that altruism done ``for the spiritual benefit'' is ``nothing but selfishness.'' Here the label is accepted.} Can I call it selfish and still get the benefit? ``Apparently I can.'' When they say ``Thank you!'', my brain feels it has done its good deed for the day. Your mileage may vary; different things restore willpower for different people, and we do not understand the deeper rules that would predict the variation.

If selfish good deeds work for you too, then ``purchase warm fuzzies and utilons separately,'' not at once, since doing both means neither is done well. If status matters to you, buy it separately too. \nb{Steven Landsburg had drawn the same line in 1997: people give because they care about the recipients or because giving feels good, and those who care about the recipients should concentrate on the worthiest charity.} A new billionaire should buy fuzzies by personally but anonymously giving a \$10,000 cashier's check to a hard-working woman about to drop out of college after her husband's hours were cut. Buy status by giving \$100,000 to the sexiest X-Prize and bragging about it for five years. Then, with cold calculation, free of scope insensitivity, ambiguity aversion, status and fuzzies, find the charity with ``the greatest expected utilons per dollar,'' and give until its marginal efficiency falls below the next one's. \nb{How far to trust such calculations is disputed. GiveWell's Holden Karnofsky argued in 2011 for preferring strong evidence of doing a lot of good to weak evidence of doing far more.}

Utilons should get at least 20 times what fuzzies get; 5\% overhead for keeping yourself altruistic seems reasonable, provided the fuzzy act really helps. Buying status is unrelated to altruism: if an X-Prize gift impresses your friends more than a speedboat of the same price, skip the speedboat, but enter it under impressing friends.

The main lesson is that all three can be bought ``far more efficiently'' separately. A \$10 million check to a breast-cancer charity, though far better than spending it on parties, will ``probably'' give far less euphoria than turning one life around in person, less to talk about than an X-Prize, and there are ``probably at least a thousand'' charities that could do ``orders of magnitude'' more good with the money. I give no figures.

If you are not a billionaire, you cannot buy in bulk, but get your fuzzies cheaply from a charity with vivid, ideally in-person beneficiaries: volunteer at a soup kitchen, or hold doors. Let your other giving validate it, but do not confuse it with buying utilons. Status is probably cheaper bought as nice clothes. For utilons, ``shut up and multiply.''
